\section{Behavioral Validation Across Models and Modalities}
\label{sec:controlled}

The preceding theory defines endpoint-relative response roles. We now test whether these roles correspond to model behavior measured independently of \decaf. We begin with the 3D Shapes dataset, whose generative factors can be changed one at a time while the remaining factors are held fixed~\citep{burgess2018understanding}. This gives exact factual--counterfactual pairs and lets us formulate three controlled learning problems in which reliance, off-path sensitivity, and effect reversal can be measured independently of \decaf. Across the controlled suite, we use \emph{ResNet-18} and a small \emph{ViT}. The base grid contains \emph{30 trained models} and \emph{180 model--factor units}, with additional dedicated checkpoints and training variants for the three behavioral tests. Comparing \decaf's decomposition with these independent behavioral measurements then tests whether each component carries the intended response semantics.

\begin{figure*}[t]
    \centering
    \begin{subfigure}[t]{0.230\linewidth}
        \centering
        \includegraphics[width=\linewidth]{figures/controlled/fig2c_e_vulnerability_scatter.pdf}
        \caption{Population-level reliance.}
        \label{fig:e-validation-population}
    \end{subfigure}\hfill
    \begin{subfigure}[t]{0.25\linewidth}
        \centering
        \includegraphics[width=\linewidth]{figures/controlled/fig2d_e_representative_trajectory.pdf}
        \caption{Within-model transition.}
        \label{fig:e-validation-trajectory}
    \end{subfigure}\hfill
    \begin{subfigure}[t]{0.234\linewidth}
        \centering
        \includegraphics[width=\linewidth]{figures/controlled/fig3b_f_fragility_scatter.pdf}
        \caption{Endpoint-null sensitivity.}
        \label{fig:f-validation}
    \end{subfigure}\hfill
    \begin{subfigure}[t]{0.234\linewidth}
        \centering
        \includegraphics[width=\linewidth]{figures/controlled/fig4c_c_vs_swap_scatter.pdf}
        \caption{Label-level reversal.}
        \label{fig:c-validation}
    \end{subfigure}
    \caption{\textbf{The three response roles track independently measured model behavior.} (a--b) Evidence follows shortcut reliance across checkpoints and through a within-model strategy transition. (c) Fragility follows off-path prediction change on endpoint-null pairs. (d) Contradiction follows pairwise label reversal. Additional intervention-separation, architecture-specific, and cross-geometry diagnostics appear in Appendix~\ref{app:controlled-results}.}
    \label{fig:controlled-validation}
\end{figure*}

\textbf{Evidence.}
We train object-shape classifiers in an environment where background wall color is correlated with the label, so models may rely on either this shortcut or the object itself. Reversing the wall--label correlation at test time gives an independent measure of shortcut reliance: a wall-dependent model loses accuracy, while a shape-dependent model remains stable. Across 52 training checkpoints, the evidence margin $E_{\mathrm{wall}}-E_{\mathrm{shape}}$ correlates $0.936$ with this reversal vulnerability, with a 90\% bootstrap interval of $[0.885,0.961]$ (Figure~\ref{fig:controlled-validation}(a)). Figure~\ref{fig:controlled-validation}(b) shows the same relation during a within-model strategy transition; complete trajectories appear in Appendix~\ref{app:controlled-results}.

\textbf{Fragility.}
We construct models for which floor color has little effect once the image is fully revealed, but different effects before full reveal. To create this difference, we expose models during training to partially revealed inputs while preserving the original shape task. In \emph{fragile training}, the model is encouraged to react to changes in floor color at these partial states. In \emph{robust training}, it is instead encouraged to remain unchanged, while \emph{neutral training} uses only the original clean task without either objective. We then measure the resulting sensitivity independently using held-out floor-color interventions. If $F$ captures endpoint-null path sensitivity, it should track this prediction-change rate. It does so closely (Figure~\ref{fig:controlled-validation}(c)); intervention separation and cross-geometry checks appear in Appendix~\ref{app:controlled-results}.

\textbf{Contradiction.}
We construct three tasks with the same object-color effect in the original wall context but different behavior when the wall context changes. In \emph{Direct}, the effect is preserved; in \emph{Gate}, it disappears; and in \emph{Invert}, it reverses. The resulting label behavior---preservation, collapse, or swap---therefore provides an independent measure of what happened to the effect. If $C$ captures opposition rather than mere weakening, it should remain near zero for Direct and Gate and increase only when the effect reverses. Across 30 models and all mismatch levels, $C$ tracks the pairwise label-swap rate with $\rho=0.961$, whereas ordinary magnitude does not ($\rho=-0.036$; Figure~\ref{fig:controlled-validation}(d)). Detailed regime separation and calibration appear in Appendix~\ref{app:controlled-results}.

\textbf{Transfer across model classes and modalities.}
We next ask whether the same response roles survive beyond controlled vision models. We use a balanced 240,000-example subset of Covertype~\citep{blackard1998covertype} with 54 natural features and append a binary context and a binary candidate factor. In one set of tasks, the candidate factor has the same effect in the reference context, while changing the context either preserves, removes, or reverses that effect. In a second set, the factor remains nearly irrelevant in the reference context but can affect predictions under the alternate context. We train \emph{135 classifiers} spanning linear, tree-based, and neural models.

Crucially, these task constructions specify what the training data attempt to induce, not what the model necessarily learns. We therefore determine preservation, inversion, and endpoint-null sensitivity directly from held-out predictions. If the response roles transfer, $E$ should track realized preservation, $C$ realized inversion, and $F$ endpoint-null prediction change. See \tablename~\ref{tab:covertype-main}. 

Remark that some models trained under the inversion construction instead suppress the candidate effect, while the fragility construction becomes ordinary endpoint evidence for logistic regression. \decaf follows the behavior realized by the trained model rather than the behavior intended by the data generator; the family-level audit appears in Appendix~\ref{app:covertype-transfer}.

\begin{table*}[!t]
    \centering
    \caption{
    \textbf{Cross-model transfer of the three response roles on Covertype.}
    Entries are Spearman correlations with held-out realized behavior across
    135 classifiers. The \decaf column uses $E$ for preservation, $C$ for
    actual inversion, and $F$ for alternate-context prediction change
    restricted to endpoint-null pairs. Brackets give 95\% joint family/seed
    bootstrap intervals. $\dagger$ SHAP interaction is available only for tree
    models. Complete baselines, model-family results, threshold-conditioned
    analyses, and measured cost appear in
    Appendix~\ref{app:covertype-transfer}.
    }
    \label{tab:covertype-main}

    \small
    \setlength{\tabcolsep}{5.0pt}

    \begin{tabular}{lllrrrr}
        \toprule
        Held-out behavior
        & Role
        & \cellcolor{decafbluehead}\decaf $\rho$
        & $\Abs$
        & $M$
        & Native SHAP
        & SHAP inter. \\
        \midrule

        Preservation
        & $E$
        & \cellcolor{decafblue}\textbf{0.864}\,[0.832,0.895]
        & 0.804
        & 0.657
        & 0.781
        & $-0.251^{\dagger}$ \\

        Actual inversion
        & $C$
        & \cellcolor{decafblue}\textbf{0.987}\,[0.973,0.997]
        & $-0.207$
        & $-0.001$
        & $-0.205$
        & $0.093^{\dagger}$ \\

        Endpoint-null change
        & $F$
        & \cellcolor{decafblue}\textbf{0.974}\,[0.942,0.988]
        & 0.588
        & 0.148
        & 0.481
        & $0.069^{\dagger}$ \\

        \bottomrule
    \end{tabular}
\end{table*}

